% Discussion
\section{Discussion}
\label{sec:discussion}

Our results provide evidence that both components of bidirectional alignment---user learning and AI responsiveness---exist in RLHF-trained conversational systems, though with weaker effects than initially hypothesized. This section interprets these findings, acknowledges limitations, and discusses implications for AI alignment research and practice.

\subsection{Interpreting the Evidence}

\subsubsection{User Learning: Small Effect, Real Signal}

The negative reformulation slope (mean = $-0.0214$, $p=0.012$, $d=-0.246$) confirms users adapt their query strategies during conversations. This learning manifests as decreasing reformulation rates---users who initially struggle to phrase queries effectively learn what works and reformulate less over subsequent turns. The small effect size (Cohen's $d$ below medium threshold) indicates learning is subtle rather than dramatic, but the statistical significance across 88 conversations suggests the signal is real.

Why is the effect small? Reformulation is a coarse-grained behavioral measure. Users may learn AI capabilities without exhibiting reformulation patterns---for instance, by asking entirely different questions based on gained understanding rather than rephrasing the same question. Our metric captures only one manifestation of learning (reformulation reduction), missing other learning signals like query sophistication increases or successful first-try queries.

The median slope = 0 finding reveals learning is not universal. Most conversations show flat slopes (no detectable learning curve), while a minority show strong negative slopes. This heterogeneity could reflect individual differences in learning aptitude, variation in task difficulty, or engagement decay confounding true learning. Future work must distinguish between learning (negative slope predicts success) and disengagement (negative slope reflects giving up) by stratifying on conversation outcomes.

\subsubsection{AI Responsiveness: Below Threshold, Above Zero}

The diversity correlation ($r=0.396$, $p<0.001$) demonstrates the RLHF-trained AI system exhibits behavioral responsiveness---response vocabulary varies with query vocabulary. However, the correlation falls below our preregistered threshold ($r>0.4$), indicating responsiveness is weaker than predicted.

We attribute this threshold failure to measurement limitations rather than absence of responsiveness. Distinct-1 captures only lexical diversity (vocabulary variation), not semantic diversity. An AI system that responds to diverse queries with semantically varied responses but overlapping vocabulary would show low Distinct-1 correlation despite genuine responsiveness. Recomputing with semantic diversity metrics (e.g., SBERT embedding variance across responses) would test whether semantic responsiveness is stronger than lexical responsiveness.

The temporal dynamics finding---correlation strengthens from $r=0.04$ (2--3 turns) to $r=0.19$ (6+ turns)---provides compelling evidence for co-adaptation. If responsiveness were a static property of the RLHF policy, correlation should be constant across conversation lengths. The observed strengthening suggests responsiveness emerges through extended interaction, consistent with alignment as a dynamic process rather than fixed policy characteristic. However, this could also reflect a confound (longer conversations accumulate more vocabulary variation), requiring partial correlation analysis to isolate true temporal effects from length artifacts.

\subsubsection{Bidirectional Alignment Framework: Partially Validated}

Our results validate the theoretical foundation of bidirectional alignment but leave the central hypothesis---coupling predicts conversation quality beyond individual metrics---untested. We have shown that both components exist (learning and responsiveness), but not that they interact to produce better conversations. This is analogous to demonstrating flour and eggs are both ingredients in bread but not yet baking bread to confirm the combination works.

The planned coupling analysis (h-m3, not executed in this study) will test whether the correlation between reformulation slope and diversity correlation within conversations predicts helpfulness ratings better than slope alone or diversity alone. If coupling adds predictive power, it validates bidirectional alignment as a meaningful construct. If not, learning and responsiveness may be independent contributors to conversation quality rather than synergistic co-adaptation.

\subsection{Limitations}

\subsubsection{L1: Small Sample for Reformulation Analysis ($n=88$)}

Our reformulation slope analysis uses only 88 conversations, representing 4.4\% of the sampled dataset (2,000 conversations with $\geq 5$ turns). This small sample results from conservative reformulation detection thresholds prioritizing precision over recall. While the sample is sufficient for statistical significance ($p=0.012$), it may not generalize to the full HH-RLHF corpus or other datasets.

\textbf{Mitigation}: Expanding to the full 161k conversation dataset and relaxing detection thresholds (at the cost of false positives) would test effect robustness. Cross-dataset validation (e.g., OpenAI API logs, Anthropic production data) would assess generalization beyond HH-RLHF.

\subsubsection{L2: Threshold Failure for Responsiveness}

The $r=0.396 < 0.4$ threshold failure weakens the claim that AI exhibits strong responsiveness. While the correlation is statistically significant and consistent with medium effect sizes in behavioral research, it falls short of the threshold we preregistered as evidence of meaningful responsiveness.

\textbf{Mitigation}: (1) Recompute with semantic diversity metrics to test whether lexical diversity underestimates effect size. (2) Apply the planned helpfulness filter (conversations with ratings $>$ median) to test whether responsiveness is stronger in high-quality conversations. (3) Revise the hypothesis to reflect weak-to-medium responsiveness ($r\approx 0.35$--0.40) rather than strong responsiveness ($r>0.4$).

\subsubsection{L3: Reformulation Detection Not Validated}

Our reformulation detection heuristic (SBERT similarity $> 0.7$ AND edit distance $> 0.3$) has not been validated against human annotations. False positives (topic shifts misclassified as reformulations) would artificially inflate slope estimates, while false negatives (reformulations missed) would underestimate slopes. The net effect on results is unknown.

\textbf{Mitigation}: Human annotation study on a 100-conversation subset would quantify precision and recall. High precision/recall ($>0.8$) would validate the heuristic; low precision/recall would require recomputing slopes with a validated detector.

\subsubsection{L4: Correlation vs Causation}

Our results show correlation (reformulation decreases over turns, response diversity correlates with query diversity) but do not establish causation. We cannot claim users' reformulation reduction \emph{causes} better alignment or that AI responsiveness \emph{causes} user satisfaction. Correlation could reflect confounds (e.g., longer conversations allow more vocabulary accumulation) or reverse causation (satisfied users reformulate less regardless of learning).

\textbf{Mitigation}: Causal intervention study where AI responsiveness is experimentally manipulated (e.g., A/B test with high-responsiveness vs low-responsiveness policies) would establish causal direction. Alternatively, instrumental variable analysis using conversation metadata as instruments could estimate causal effects from observational data.

\subsubsection{L5: Coupling Hypothesis Untested}

The core hypothesis---bidirectional alignment coupling predicts helpfulness---remains untested. Our results establish that components exist but not that they interact productively. This leaves the main theoretical claim unvalidated.

\textbf{Mitigation}: Complete the planned h-m3 coupling analysis: regress helpfulness ratings on coupling strength (correlation between slope and diversity per conversation) while controlling for slope alone and diversity alone. Significant coupling coefficient would validate the interaction hypothesis.

\subsection{Broader Impact}

\subsubsection{Implications for Alignment Measurement}

Our work suggests alignment should be evaluated not just through static policy quality (preference accuracy, single-turn helpfulness) but through conversation dynamics. Two AI systems with identical per-turn response quality may differ in bidirectional alignment---one elicits user learning and exhibits responsiveness, while the other does not. Standard RLHF evaluation would rate them equally; bidirectional metrics would differentiate them.

This has practical implications for deployed systems. Monitoring reformulation slopes and diversity correlations in production logs could detect alignment degradation that aggregate helpfulness ratings miss. For instance, if users increasingly reformulate queries without learning (flat slopes despite high reformulation rates), this signals the AI is not adapting effectively despite potentially high per-turn quality.

\subsubsection{Implications for Training Methods}

Current RLHF optimizes AI policy to maximize human preference predictions, treating users as static evaluators. Our results suggest users are not static---they learn during conversations. Future alignment training could incorporate bidirectional objectives:

\begin{itemize}
\item \textbf{User learning objective}: Encourage responses that reduce future reformulation rates (teaching users what queries work).
\item \textbf{AI responsiveness objective}: Train policies that adjust output characteristics (diversity, detail, formality) based on query patterns.
\item \textbf{Coupling objective}: Maximize correlation between user learning rate and AI responsiveness within conversations, directly optimizing for co-adaptation.
\end{itemize}

These objectives would shift RLHF from unidirectional preference matching to bidirectional co-adaptation, potentially producing AI systems that are more helpful not just in individual responses but across extended interactions.

\subsubsection{Societal Considerations}

Optimizing for bidirectional alignment raises concerns about user dependency. If AI systems are trained to be maximally responsive to user patterns, users may become reliant on AI communication styles and struggle with less adaptive systems. This is analogous to autocomplete making users worse at spelling---convenience during use may reduce transferable skills.

We recommend balancing bidirectional alignment optimization with user autonomy. AI systems should encourage learning (reducing reformulation through informative responses) without creating dependency (over-responsive AI that accepts any query formulation reduces user incentive to learn effective communication). Monitoring both coupling strength (alignment quality) and user query diversity over time (skill maintenance) would detect unhealthy dependency patterns.

\subsection{Future Directions}

Our partial validation of bidirectional alignment components opens several research directions:

\begin{enumerate}
\item \textbf{Complete coupling hypothesis test}: Execute h-m3 to test whether correlation between slope and diversity predicts helpfulness beyond individual metrics.

\item \textbf{Semantic diversity reanalysis}: Recompute responsiveness with SBERT embedding variance instead of Distinct-1 to test whether semantic responsiveness exceeds lexical responsiveness.

\item \textbf{Success stratification}: Test whether negative reformulation slopes predict task success (learning interpretation) or failure (disengagement interpretation).

\item \textbf{Causal intervention}: A/B test with manipulated AI responsiveness levels to establish causal direction between responsiveness and user satisfaction.

\item \textbf{Cross-dataset validation}: Replicate findings on non-HH-RLHF datasets (e.g., ShareGPT, Anthropic production logs) to assess generalization.

\item \textbf{Bidirectional RLHF training}: Implement training algorithms that optimize for coupling strength rather than static preference predictions.
\end{enumerate}

The bidirectional alignment framework provides a lens for understanding AI-human interaction that complements existing unidirectional evaluation methods. By measuring how both participants adapt during conversations, we can design AI systems that are not just helpful in isolation but genuinely aligned through interaction.
